\section{Empirical Evaluation}
\label{section:evaluation}
% Please concretely describe how you plan to evaluate the effectiveness of your method by completing the following parts.

\subsection{Performance Metrics}
% Please list a few key performance metrics that you plan to use in your evaluation.
% Typical examples in deep learning settings include accuracy, precision, recall, F1 score, AUROC, mean average precision (mAP), Intersection over Union (IoU), mean squared error (MSE), mean absolute error (MAE), or other task-specific evaluation metrics.
% If you plan to design a new evaluation metric, please also explain why it is necessary and why it provides a reasonable measure of model performance.
For our Image-to-Sequence OMR task, we plan to use the following performance metrics:
\begin{itemize}
    \item \textbf{Symbol Error Rate (SER):} Similar to Character Error Rate (CER) in traditional OCR tasks, SER calculates the Levenshtein distance (accounting for insertions, deletions, and substitutions) between the predicted Jianpu token sequence and the ground truth. This is the most critical metric for evaluating the quality of sequence generation.
    \item \textbf{Sequence-level Accuracy (Exact Match Rate):} This metric calculates the percentage of completely correct musical measures or lines. It is a strict metric that ensures the entire musical context within a given window is perfectly translated without any token errors.
\end{itemize}

\subsection{Baseline Methods}
% Please provide a list of the necessary baseline methods for comparison. For example:
% \begin{itemize}
% \item Baseline 1 [authors, year]: A standard supervised learning model trained on the chosen dataset, serving as a strong baseline.
% \item Baseline 2 [authors, year]: A commonly used deep neural network architecture for the target task.
% \item Baseline 3 [authors, year]: A recent state-of-the-art or representative method closely related to your proposed approach.
% \item $\cdots$
% \end{itemize}
% In the above, please use ``\textbackslash citep{...}'' to add the references.
Since end-to-end models specifically designed for Staff-to-Jianpu translation are extremely rare, we will compare our proposed architecture with the following baselines:
\begin{itemize}
    \item \textbf{Pipeline-based SOTA \citep{oemer2023}:} Oemer, an open-source OMR system utilizing U-Net for semantic segmentation. We will evaluate its intermediate ability to extract staff symbols, which would hypothetically require an additional heuristic converter to produce Jianpu.
    \item \textbf{Classic Image-to-Sequence \citep{shi2015crnn}:} CRNN (Convolutional Recurrent Neural Network) trained with CTC loss. This is a standard and robust baseline for visual sequence recognition tasks, which we will adapt for our custom Jianpu token vocabulary.
    \item \textbf{Modern Vision-Language Model \citep{li2021trocr}:} TrOCR (Transformer-based Optical Character Recognition). As a pre-trained Vision-Encoder-Decoder model, it represents a strong architectural baseline for state-of-the-art end-to-end sequence generation.
\end{itemize}
As a first step, we additionally run an internal \emph{encoder ablation} within our own Vision-Encoder-Decoder (a pretrained Vision Transformer vs.\ a from-scratch ResNet backbone) to isolate the contribution of the visual representation; preliminary results are reported in Section~\ref{section:exp}. The three external baselines above are then evaluated against the resulting model in the final report.

\subsection{Benchmark Tasks or Datasets}
% Please provide a list of the benchmark tasks or datasets that you plan to use in your evaluation. For example:
% \begin{itemize}
% \item CIFAR-10, CIFAR-100, or ImageNet for image classification.
% \item COCO or Pascal VOC for object detection and segmentation.
% \item LibriSpeech for speech recognition.
% \item UCI, M4, or other standard datasets for time-series forecasting.
% \item MovieLens or Amazon Review datasets for recommendation.
% \item Multi-modal benchmarks such as VQA or Flickr30k for vision-language tasks.
% \end{itemize}

% If your project belongs to the Dataset and Benchmark Track, please describe in detail how you plan to construct your dataset or benchmark, including the data collection pipeline, annotation process, quality control procedure, and evaluation protocol.
Because there is no standard public dataset for Staff-to-Jianpu translation, we will evaluate our method on the following custom data setups:
\begin{itemize}
    \item \textbf{Custom Synthetic Dataset:} A large-scale training and validation set generated programmatically. We will render digital score files (e.g., MusicXML or MIDI) into both standard Staff images and their corresponding Jianpu token sequences using rendering engines (such as MuseScore or LilyPond) to provide massive paired data.
    \item \textbf{Real-World Music Test Set:} A manually collected and annotated test set consisting of scanned or photographed printed sheet music commonly used in music ensembles. This serves as the ultimate benchmark to evaluate our model's generalization capability and practical utility in real-world scenarios.
\end{itemize}
